% Solutions of Book 4, chapter 14 (files-pipes-signals).
\section*{Chapter \ref{ch:b2:files-pipes-signals} --- Files, Pipes and Signals}

\begin{solution}{exo:b2:files-pipes-signals:1}
The calls return 10, 10, 6 and 0: the offset goes 0, 10, 20, 26, and the fourth read finds nothing left. The third
read copies \code{uvwxyz} into the first six bytes of \code{buf} and leaves the other four as the second read left
them, so \code{buf} holds \code{uvwxyzqrst}: only the first 6 bytes, the count \code{read} returned, are data. After
\code{lseek(fd, 20, SEEK_SET)} the offset starts at 20: one read returns 6, \code{uvwxyz}, and the next 0.
\end{solution}

\begin{solution}{exo:b2:files-pipes-signals:2}
The child duplicates \code{fd[0]} onto 0 and closes both \code{fd[0]} and \code{fd[1]} before its \omterm{def:b2:processes-and-system-calls:exec}{exec}, so
\code{wc} holds no write end. The parent closes \code{fd[0]} but keeps \code{fd[1]}: it is the only holder, and by
\cref{prop:b2:files-pipes-signals:eof} \code{wc} waits as long as it is open. The missing \code{close(fd[1])} goes
right after the \code{write}, before \code{waitpid}. After \code{waitpid} it is too late: the parent waits for
\code{wc}, which waits for end of file, which waits for the parent's \code{close}. Each waits for the other, forever.
\end{solution}

\begin{solution}{exo:b2:files-pipes-signals:3}
\omcode{code/bachelor-2/14-files-pipes-signals/exercises/mycp.c}{24}{44}{Copying with the source's permission
bits}
\code{fstat} reads the source's \omterm{def:b2:files-pipes-signals:inode}{inode}, so the mode is the source's. \code{O_EXCL} with \code{O_CREAT} makes
\code{open} fail if the destination exists, in the same \omterm{def:b2:processes-and-system-calls:syscall}{system call} that would create it, so no other program can
create it in between. The mode passed to \code{open} is reduced by the \omterm{def:b2:processes-and-system-calls:process}{process}'s umask, which usually clears the
write bits of the group and others, so \code{fchmod} sets the bits exactly at the end. The inner loop writes again
until the \code{n} bytes read are all written, since a \code{write} may write fewer. The test copies 70\,000 bytes
of mode \code{751} with $b = 1$, 7, 4096 and 100\,000, and checks that a second copy to the same name fails.
\end{solution}

\begin{solution}{exo:b2:files-pipes-signals:4}
One \code{write} for each read that returned bytes: $\lceil n / b \rceil$. With the $\lceil n/b \rceil + 1$ reads of
\cref{prop:b2:files-pipes-signals:calls}, two \code{open}s and two \code{close}s, the total is
$2\lceil n/b \rceil + 5$ (the loader's own calls before \code{main}, about thirty, come on top). For $n = 10^6$, the
total is under 300 when $\lceil 10^6 / b \rceil \le 147$, that is $b \ge 10^6 / 147$, so $b \ge 6803$:
$\lceil 10^6 / 6803 \rceil = 147$ gives 299 calls, while $b = 6802$ gives 148 writes and 301.
\end{solution}

\begin{solution}{exo:b2:files-pipes-signals:5}
\omoutput{code/bachelor-2/14-files-pipes-signals/outputs/exo_nonewline.txt}{Two lines, both with the
\code{A}.}
Into a file, the output is fully buffered: \code{A} waits in the buffer, \code{fork} copies it, each \omterm{def:b2:processes-and-system-calls:process}{process} appends
\code{B} and a newline, and each \omterm{def:b2:files-pipes-signals:buffering}{flushes} \code{AB} at exit. On a terminal the output is line buffered, but a line is
flushed at its newline and \code{A} has none, so it is still in the buffer at the fork: the terminal shows \code{AB}
twice as well. The newline matters, not the terminal. \code{fflush(stdout)} before \code{fork} prints \code{A},
then \code{B} twice.
\end{solution}

\begin{solution}{exo:b2:files-pipes-signals:6}
\code{dup2(1, fd)} makes \code{fd} point to the description of descriptor 1, the terminal or the capture file, and
closes what \code{fd} pointed to: \code{out.txt}'s only description, so the file stays as \code{open} left it,
created and truncated. \code{close(fd)} then removes the extra copy of descriptor 1, and \code{echo} writes to 1 as
always: the text goes to the parent's output, and \code{out.txt} has 0 bytes. Read \code{dup2(old, new)} as ``make
\code{new} a copy of \code{old}''.
\end{solution}

\begin{solution}{exo:b2:files-pipes-signals:7}
\omcode{code/bachelor-2/14-files-pipes-signals/exercises/tee.c}{23}{46}{\code{tee}: one read, several
writes}
Descriptor 1 is placed first in the table of outputs, so the loop treats it like the files. \code{O_APPEND} makes the
kernel move the offset to the end before each write, which is what \code{-a} asks even if another program writes the
file meanwhile. A file that cannot be opened is reported and skipped; the others are still written, and the status
says something failed. \code{write_all} loops as in \code{mycp}. The test sends 18\,000 bytes through two files,
then appends a line with \code{-a}.
\end{solution}

\begin{solution}{exo:b2:files-pipes-signals:8}
Suppose the \omterm{def:b2:files-pipes-signals:pipe}{pipe} is empty. By the definition of a \omterm{def:b2:files-pipes-signals:pipe}{pipe}, a read on it returns 0 when no write end is open anywhere,
which is the ``if'' direction. Conversely, while some \omterm{def:b2:processes-and-system-calls:process}{process} holds an open descriptor of the write end, that \omterm{def:b2:processes-and-system-calls:process}{process}
may still write: the kernel cannot know that no byte will come, so the read waits, and returns 0 only once the last
such descriptor is closed (by \code{close}, by \code{dup2} over it, or by the \omterm{def:b2:processes-and-system-calls:process}{process}'s end, which closes all its
descriptors). Closing a descriptor of the read end changes nothing, and neither does a reader's copy of the write end
being ``its own'': it counts like any other.

In the method, \omterm{def:b2:files-pipes-signals:pipe}{pipe} $i$'s write end is created in the parent, inherited by child $i$, and closed by the parent right
after that fork, before child $i + 1$ exists; child $i$ moves it to descriptor 1 and closes the original. So the only
descriptor of \omterm{def:b2:files-pipes-signals:pipe}{pipe} $i$'s write end is child $i$'s descriptor 1, and child $i + 1$ sees end of file exactly when child
$i$ has ended (or closed its output). By induction on $i$: child 1 ends at some moment; child 2 sees end of file only
after that, and, reading to the end before exiting, ends only after that; and so on to child $k$, which sees end of
file only after children $1, \dots, k-1$ have all ended.
\end{solution}

\begin{solution}{exo:b2:files-pipes-signals:9}
\omcode{code/bachelor-2/14-files-pipes-signals/exercises/capture_cmd.c}{14}{28}{Two children: the command, and one
that feeds it}
The function grows its result while it reads. The naive version hangs because two \omterm{def:b2:files-pipes-signals:pipe}{pipes} are full at once. \code{cat}
copies its input to its output; nobody reads the output yet, so after 64 KiB that \omterm{def:b2:files-pipes-signals:pipe}{pipe} is full and \code{cat} waits in
\code{write}, and stops reading. Then the input \omterm{def:b2:files-pipes-signals:pipe}{pipe} fills with the next 64 KiB, and the parent waits in its own
\code{write}. Neither moves again. With 300\,000 bytes, more than the two \omterm{def:b2:files-pipes-signals:pipe}{pipes} and \code{cat}'s buffer hold, it
happens every time; with a few KiB, never. The cure is to write and read at the same time: here a second child writes
the input while the parent reads the output. The parent must close its copy of \code{to[1]}, or the command never sees
end of file; the feeder ends after its write, closing the last copy. The test checks \code{cat}, a status of 3, and
the 300\,000 bytes.
\end{solution}

\begin{solution}{exo:b2:files-pipes-signals:10}
\omcode{code/bachelor-2/14-files-pipes-signals/exercises/minish2.c}{25}{45}{Applying and removing the
redirections, in the child}
\omcode{code/bachelor-2/14-files-pipes-signals/exercises/minish2.c}{57}{79}{Cutting the line into stages and wiring
them}
The line is cut at each \code{|} into stages, wired by \cref{met:b2:files-pipes-signals:pipeline}. Each child first
installs its \omterm{def:b2:files-pipes-signals:pipe}{pipe} ends, then applies its redirections, which remove themselves from the words before the \omterm{def:b2:processes-and-system-calls:exec}{exec}. In
that order a \code{>} replaces the \omterm{def:b2:files-pipes-signals:pipe}{pipe} on descriptor 1, so \code{a > f | b} writes \code{a}'s output to \code{f},
and \code{b} reads an empty \omterm{def:b2:files-pipes-signals:pipe}{pipe}: end of file as soon as \code{a} ends. That is also what a real shell does. In the
other order the \omterm{def:b2:files-pipes-signals:pipe}{pipe} would overwrite the redirection. The test runs a redirected \code{sort}, a three-stage pipeline,
and \code{false | true}.
\end{solution}

\begin{solution}{exo:b2:files-pipes-signals:11}
\omcode{code/bachelor-2/14-files-pipes-signals/listings/fit_copy.py}{11}{26}{Least squares on relative
errors}
\omoutput{code/bachelor-2/14-files-pipes-signals/outputs/fit_copy.txt}{The fit, on one run of the author's
laptop.}
The model is linear in $a$ and $c$; the script uses the read counts of the measurement and minimises the relative
errors, so that the fast rows count as much as the slow ones. On this run $a \approx 375$~ns for one read and its
write, and $c \approx 0.2$~ns per byte; the model is within about 10~\% of every row. A read and a write are two
\omterm{def:b2:processes-and-system-calls:syscall}{system calls}, about 190~ns each: more than the 70~ns of the empty call of \cref{ch:b2:processes-and-system-calls},
because each finds the file, checks the offset and copies data. The two terms are equal when $a\,n/b = c\,n$, that is
$b = a / c \approx 375 / 0.215 \approx 1700$ bytes. Beyond about 2 KiB the bytes cost more than the calls, which is
where the curve of the chapter flattens.
\end{solution}

\begin{solution}{exo:b2:files-pipes-signals:12}
\omcode{code/bachelor-2/14-files-pipes-signals/exercises/minish2.c}{101}{128}{The prompt reports; the handler
reaps}
The handler is that of \code{sigchld_reap.c}: it reaps every ended child with \code{WNOHANG} and counts them, and the
main loop, never the handler, prints the report before the next prompt. A background line is wired as any other and
not waited for. The danger is the handler's \code{waitpid(-1, ...)} reaping a foreground child before the shell's
own \code{waitpid}, which would then fail and lose the status. Two measures prevent it: \code{SIGCHLD} is blocked from
before the forks until the foreground stages are reaped, and the shell waits for its own PIDs only, never for any
child. When the mask is restored, the pending \omterm{def:b2:files-pipes-signals:signal}{signal} runs the handler, which finds only background processes left. A
child inherits the blocked mask, and an \omterm{def:b2:processes-and-system-calls:exec}{exec} keeps it, so each child unblocks \code{SIGCHLD} before its \omterm{def:b2:processes-and-system-calls:exec}{exec}. The
test runs \code{sleep 0.2 &}, then \code{sleep 0.6}, and expects one report.
\end{solution}

\begin{solution}{pb:b2:files-pipes-signals:1}
\textbf{1.} 3: \code{open} returns the lowest free descriptor, and 0, 1 and 2 are the standard streams, open since
the program started.

\textbf{2.} The \omterm{def:b2:files-pipes-signals:descriptor}{open file description} of \code{words.txt}, at offset 0; 3 is free again. Every child inherits it,
but children 2 to 6 put a \omterm{def:b2:files-pipes-signals:pipe}{pipe} over their descriptor 0, so only the first, \code{tr -cs}, reads it.

\textbf{3.} \code{close} freed 3, which is again the lowest free descriptor.

\textbf{4.} The fork copies the \omterm{def:b2:files-pipes-signals:descriptor}{descriptor table}, not the description: parent and child share one offset. The child's
11 bytes move it to 11, and the parent, which waited, writes its 12 bytes from there: 23 bytes, child first.

\textbf{5.} The child would have a description of its own, at offset 0, and the parent's would stay at 0 too. The
child writes its 11 bytes at 0, then the parent writes its 12 bytes at 0 over them: the file holds
\code{parent line} and its newline, 12 bytes.

\textbf{6.} The type and the permission bits (\code{-rw-r--r--}, or \code{l} for the link), the \omterm{def:b2:files-pipes-signals:inode}{link count}, the \omterm{def:b2:rust-ownership:ownership}{owner}
and group, the size, and the time of the last change. The name comes from the directory, and the
\code{-> hours.txt} from the symbolic link's contents.

\textbf{7.} Child 1, \code{tr -cs}: 0 is \code{words.txt}, 1 the write end of \omterm{def:b2:files-pipes-signals:pipe}{pipe} 1. Child $i$ for $i = 2$ to 5: 0
is the read end of \omterm{def:b2:files-pipes-signals:pipe}{pipe} $i - 1$, 1 the write end of \omterm{def:b2:files-pipes-signals:pipe}{pipe} $i$. Child 6, \code{head}: 0 is the read end of \omterm{def:b2:files-pipes-signals:pipe}{pipe} 5, and 1
is the parent's standard output.

\textbf{8.} Three. At the start of iteration $i$ the parent holds one, \code{prev}; \code{pipe} adds two, the fork
copies them all to the child, and then the parent closes \code{prev} and \code{fd[1]}, keeping only \code{fd[0]}.
So it holds three from the \code{pipe} to the closes, and one between iterations; after the last iteration, none.

\textbf{9.} \code{tr -cs}, reading a regular file: its read returns 0 after the 440 bytes. Among the \omterm{def:b2:files-pipes-signals:pipe}{pipe} readers,
child 2 is first: its end of file comes when child 1 ends, since child 1 holds the only write end of \omterm{def:b2:files-pipes-signals:pipe}{pipe} 1.

\textbf{10.} The last line of the input may be the first line of the output, so \code{sort} can write nothing before
end of file. \code{tr} translates each byte on its own and writes as soon as its buffer is full.

\textbf{11.} None: end of file depends only on write ends. But if \code{head} stopped early with more data on the
way, the stage before it would not die of \code{SIGPIPE}, since a read end would still be open in the parent. It
would fill the \omterm{def:b2:files-pipes-signals:pipe}{pipe} and wait in \code{write} forever. With this small text no stage blocks, but in general one would.

\textbf{12.} \code{false | true} has status 0, and \code{grep word missing.txt | sort} has status 0 though the file
does not exist. A shell can collect every stage's status, as \code{run_pipeline}'s loop does before it keeps only the
last, and return the first that is not 0. Some shells offer this as an option (\code{pipefail}).

\textbf{13.} With $n = 100\,000$: $b = 1$, $100\,000 + 1$; $b = 10$, $10\,000 + 1$; $b = 100$, $1000 + 1$;
$b = 1000$, $100 + 1$; $b = 4096$, $\lceil 24.4 \rceil + 1 = 26$; $b = 10\,000$, $10 + 1$; $b = 65\,536$,
$\lceil 1.53 \rceil + 1 = 3$. Each count matches.

\textbf{14.} $b = 3$: $\lceil 33\,333.3 \rceil + 1 = 33\,335$. $b = 99\,999$: $2 + 1 = 3$. $b = 100\,000$:
$1 + 1 = 2$. $b = 100\,001$: also 2.

\textbf{15.} On the run printed, $1619.3 / 380.2 \approx 4.3$, while the read counts give
$4\,194\,305 / 1\,048\,577 \approx 4.0$. At small $b$ the time is almost exactly the number of \omterm{def:b2:processes-and-system-calls:syscall}{system calls} times a
constant.

\textbf{16.} Past a few KiB the calls are few: 1025 reads and writes at $b = 4096$, about 0.4~ms at the 375~ns of
\cref{exo:b2:files-pipes-signals:11}. The time left is the moving of the 4 MiB from the kernel to the buffer and back,
the $c\,n$ term, which no buffer size reduces.

\textbf{17.} By \cref{prop:b2:files-pipes-signals:calls} with $n = 440$ and $b = 4096$: $\lceil 0.107 \rceil + 1 = 2$
reads, one returning the 440 bytes and one returning 0.

\textbf{18.} \code{putc} is a function call that stores a byte in the library's buffer, a few nanoseconds; the \omterm{def:b2:processes-and-system-calls:syscall}{system
call} comes once per 4096 bytes, 1024 times for 4 MiB. That is the same number of writes as the copy with
$b = 4096$; what remains is four million calls to \code{putc}, which the raw copy does not make.

\textbf{19.} \omterm{def:b2:files-pipes-signals:pipe}{Pipe} 5's read end is closed in every \omterm{def:b2:processes-and-system-calls:process}{process} (\code{head} has exited, and the parent closed its copy),
so the write raises \code{SIGPIPE} in \code{sort}.

\textbf{20.} By default \code{SIGPIPE} kills the writer: \code{sort} ends by \omterm{def:b2:files-pipes-signals:signal}{signal} 13, which \code{run_pipeline}
records as $128 + 13 = 141$, then ignores, since only \code{head}'s status is the pipeline's.

\textbf{21.} Death by \code{SIGPIPE} is the normal end of a writer whose reader has stopped wanting more, as with
\code{yes | head -3}; a shell that reported it would print a message under most pipelines ending in \code{head}. And
the dead stage's status is not the pipeline's anyway.

\textbf{22.} Every later \code{write} fails at once with \code{EPIPE}. Not checking, the program goes on computing
and writing into nothing, until its work runs out; a program like \code{yes} would never stop.

\textbf{23.} No. \code{sort -rn} writes only after its input has ended, and its output, a few hundred bytes, leaves
its buffer in one \code{write} at exit. The \omterm{def:b2:files-pipes-signals:pipe}{pipe} holds 65\,536 bytes, so that write succeeds at once, whether or not
\code{head} has yet read anything; \code{head} has not exited before it, since it waits for its first read. No write
happens after \code{head}'s exit.

\textbf{24.} The buffer-size law: a copy of $n$ bytes with a buffer of $b$ bytes makes $\lceil n/b \rceil + 1$
reads. For $n = 100\,000$ it predicts, from $b = 1$ to $b = 65\,536$, 100\,001, 10\,001, 1001, 101, 26, 11 and 3
reads, the seven counts of \cref{out:b2:files-pipes-signals:counts}.
\end{solution}
